\documentclass{article}
\usepackage{amsmath,amssymb,hyperref}
\title{Product-indexed directed parallel and native lax-cone objects}
\date{15 November 2026}
\begin{document}\maketitle
\paragraph{Status.} Written construction independently reviewed in the stated product-index scope. This derives a binary assembly identification from the original directed parallel, rather than assuming parallel already acts on every newly admitted residual right adjoint. The conclusion is product-indexed, not arbitrary shared-index synchronization.

\paragraph{Ambient.} Let $Q$ be locally cocomplete with separately cocontinuous composition. Let $T:Q\times Q\to Q$ be a specified normal lax parallel operation, separately locally cocontinuous. Its directed compositor has orientation
\[
 T(v_1,v_2)\,T(u_1,u_2)\longrightarrow T(v_1u_1,v_2u_2).
\]
Use the paired enriched-category/module extension $T_*$. An enriched category pair has object pairs and hom $T(A(a',a),B(b',b))$; module pairs have the analogous components. Actions and the module comparison use the directed compositor followed by the module quotient. No invertibility of this comparison is needed below. All object, arrow and coproduct indices are small in the chosen universes.

\paragraph{First extension: normal lax diagrams.} For a normal lax diagram $F:D\to\operatorname{Mod}(Q)$ from a small ordinary category, form $A_F$ with objects $(d,x)$ and homs
\[
 A_F((d',x'),(d,x))=\coprod_{a:d\to d'}F(a)(x',x).
\]
Composition uses the quotient into $F(b)\odot F(a)$ and the possibly noninvertible diagram compositor into $F(ba)$. The lax associativity equation proves associativity; normality supplies the identity summand and the unit equations. The old argument never needed to invert a nonidentity diagram compositor. Modules $W:X\to A_F$ are therefore exactly columns $W_d:X\to Fd$ with action cells $F(a)\odot W_d\to W_{d'}$ obeying the lax cone equations, and module maps are their modifications. This is a universal property for cones to the specified lax diagram, not a claim that the diagram has become pseudo.

\paragraph{Binary theorem.} Let $F:D\to\operatorname{Mod}(Q)$ and $G:E\to\operatorname{Mod}(Q)$ be normal lax diagrams. Their paired diagram
\[
 H:D\times E\to\operatorname{Mod}(Q),\qquad
 H(d,e)=T_*(Fd,Ge),\quad H(a,b)=T_*(F(a),G(b))
\]
has compositor obtained first from that of $T_*$ and then those of $F,G$. There is a canonical isomorphism of enriched categories, up to the specified extent coherence,
\[
 T_*(A_F,A_G)\ \cong\ A_H.
\]
It sends $((d,x),(e,y))$ to $((d,e),(x,y))$. On homs the map is precisely separate coproduct preservation:
\[
 T\!\left(\coprod_aF(a)(x',x),\coprod_bG(b)(y',y)\right)
 \cong\coprod_{a,b}T(F(a)(x',x),G(b)(y',y)).
\]

\paragraph{Proof of the multiplication assertion.} Precompose both candidate multiplication maps with the coproduct injections for $a,a',b,b'$ and with the quotient maps defining the two module composites. On these generators, both sides first use the same directed compositor of $T$, then the diagram composites for $F$ and $G$. Naturality and associativity of that compositor identify the pastes, including the exterior module actions. Quotient and coproduct universality then identify the full multiplication maps. Identity summands match by normality. The inverse hom maps are the coproduct isomorphisms, not inverses of effectful interchangers. This proves an enriched-category isomorphism, hence an equivalence of the corresponding incoming module hom categories and lax-cone universal properties.

\paragraph{Projection and native cone compatibility.} The universal projection $p_j:A_F\to Fj$ has components
\[
 p_j(y,(i,x))=\coprod_{a:i\to j}F(a)(y,x).
\]
Its actions are the same quotient/compositor maps used in $A_F$. For an arrow $d:j\to j'$ the projection cone cell $F(d)\odot p_j\to p_{j'}$ uses the summand composite $F(d)\odot F(a)\to F(da)$. Consequently separate coproduct preservation identifies $T_*(p_j,p_k)$ with $p_{(j,k)}$ under the displayed enriched-category comparison. The generating cone cells agree by the same lax-compositor calculation. This supplies compatibility with the actual universal cone, not merely an isomorphism between unmarked enriched categories. No arbitrary diagram-morphism naturality is asserted here: transformations with module components may induce module comparisons, not enriched functors.

\paragraph{Higher assembly coherence.} If the given parallel has coherent associativity and unit comparisons, the corresponding product-indexed identifications inherit them: on objects they are the usual reassociations of finite tuples; on homs they are the canonical maps from independently indexed coproducts and the original parallel comparisons. Checking after coproduct injections reduces the coherence diagrams to those already required of $T$. This assertion inherits whatever strength the original comparisons have; it does not create symmetry, invertible interchange, or a trace law.

\paragraph{What shared synchronization adds.} If both processes use the same index $D$, independent assembly has index $D\times D$. Shared assembly instead uses $H\Delta$, where $\Delta:D\to D\times D$ selects equal index positions and equal diagram arrows. The two constructions have different objects and hom summands. The theorem above does not identify $A_H$ with $A_{H\Delta}$. Their comparison is the separate synchronization/restriction problem for which selected adjoints and Beck--Chevalley data were developed. This isolates the actual extra obligation instead of conflating parallel with categorical product. It also explains why the earlier Prof disjoint-union product counterexample does not contradict this product-indexed theorem.

\paragraph{Relative and feedback scope.} Start with the old relative context/content presentation and its independently derived normal lax parallel, then use its separate homwise cocontinuous extension as $T$. The theorem is about that derived operation on enriched cone presentations. It does not declare all new objects or modules pure, nor assert total parallel on the lax-action right-adjoint enlargement. Iterated shared restriction, feedback and evaluation of such presentations still require their own compatibility proof. The machinery is enriched/module and coproduct calculus; no novelty claim is made.
\end{document}
